\documentclass[../DoAn.tex]{subfiles}
\begin{document}

\section{Giải pháp multi-tenant RAG chatbot cho tư vấn bán hàng nội thất}
\label{section:5.1}

\subsection{Bài toán đặt ra}
\label{subsection:5.1.1}

Bài toán trung tâm của đồ án là xây dựng một hệ thống chatbot hỗ trợ tư vấn bán hàng nội thất cho nhiều cửa hàng trên cùng một nền tảng. Điểm khó của bài toán không chỉ nằm ở việc xây dựng một chatbot có khả năng trả lời bằng ngôn ngữ tự nhiên, mà còn ở việc tổ chức hệ thống sao cho mỗi cửa hàng có dữ liệu, chatbot, hội thoại và yêu cầu mua hàng riêng. Nếu hệ thống chỉ có một chatbot dùng chung cho tất cả cửa hàng, chatbot có thể trả lời lẫn dữ liệu giữa các cửa hàng, gây sai lệch trong quá trình tư vấn. Ngược lại, nếu xây dựng một hệ thống riêng cho từng cửa hàng, việc triển khai, vận hành và bảo trì sẽ trở nên phức tạp khi số lượng cửa hàng tăng lên.

Bên cạnh đó, dữ liệu sản phẩm nội thất có tính thay đổi theo từng cửa hàng. Mỗi cửa hàng có danh mục sản phẩm, cách mô tả, mức giá, chính sách bảo hành, chính sách vận chuyển và phong cách tư vấn riêng. Vì vậy, chatbot không thể chỉ dựa vào kiến thức chung của mô hình ngôn ngữ để trả lời. Chatbot cần có khả năng truy xuất thông tin từ dữ liệu phù hợp với cửa hàng đang được tư vấn. Đây là lý do hệ thống cần kết hợp mô hình multi-tenant với hướng tiếp cận Retrieval-Augmented Generation. Multi-tenant giúp tổ chức nhiều cửa hàng trên cùng một nền tảng nhưng vẫn tách biệt dữ liệu, còn RAG giúp chatbot trả lời dựa trên dữ liệu được truy xuất từ knowledge base.

Một vấn đề khác là quá trình tư vấn bán hàng không kết thúc ở việc chatbot trả lời câu hỏi. Trong thực tế, khi khách hàng quan tâm đến sản phẩm hoặc có ý định mua hàng, cửa hàng cần tiếp nhận thông tin đó để tiếp tục xử lý. Nếu chatbot chỉ dừng lại ở trả lời hội thoại, hệ thống sẽ thiếu liên kết với nghiệp vụ bán hàng phía sau. Vì vậy, đồ án cần thiết kế một luồng xử lý trong đó hội thoại có thể phát sinh lead hoặc yêu cầu mua hàng, giúp cửa hàng theo dõi và xử lý tiếp sau khi khách hàng đã trao đổi với chatbot.

Từ các vấn đề trên, đóng góp chính của giải pháp là kết hợp ba yếu tố: mô hình multi-tenant, chatbot RAG và luồng nghiệp vụ bán hàng. Giải pháp hướng đến việc giải quyết đồng thời ba yêu cầu: tư vấn dựa trên dữ liệu phù hợp, tách biệt dữ liệu giữa các tenant, và chuyển kết quả hội thoại thành thông tin nghiệp vụ để cửa hàng tiếp tục xử lý.

\subsection{Giải pháp đề xuất}
\label{subsection:5.1.2}

Giải pháp được đề xuất là tổ chức hệ thống theo mô hình multi-tenant RAG chatbot. Trong mô hình này, mỗi cửa hàng được biểu diễn như một tenant trong hệ thống. Các thực thể quan trọng như chatbot, nguồn dữ liệu knowledge base, hội thoại, lead và yêu cầu mua hàng đều được gắn với tenant tương ứng. Nhờ đó, hệ thống có thể phục vụ nhiều cửa hàng trên cùng một nền tảng, nhưng vẫn đảm bảo dữ liệu của từng cửa hàng được quản lý riêng biệt.

Ở tầng nghiệp vụ, hệ thống cung cấp các chức năng quản lý tenant, quản lý thành viên cửa hàng, quản lý chatbot, quản lý nguồn dữ liệu knowledge base, quản lý hội thoại, quản lý lead và quản lý yêu cầu mua hàng. Quản trị hệ thống có vai trò tạo và quản lý các tenant, trong khi quản trị cửa hàng chịu trách nhiệm cấu hình chatbot và nguồn dữ liệu của cửa hàng mình. Nhân viên cửa hàng theo dõi hội thoại, lead và yêu cầu mua hàng phát sinh từ quá trình tư vấn. Cách phân chia này giúp hệ thống phản ánh đúng các vai trò trong bài toán thực tế, đồng thời hỗ trợ kiểm soát quyền truy cập theo tenant.

Ở tầng xử lý hội thoại, khi người dùng gửi tin nhắn đến chatbot của một cửa hàng, hệ thống trước hết xác định chatbot và tenant tương ứng. Sau đó, yêu cầu hội thoại được chuyển đến AI service cùng với thông tin ngữ cảnh cần thiết. AI service thực hiện truy xuất dữ liệu liên quan từ knowledge base của tenant đó, rồi sử dụng kết quả truy xuất làm ngữ cảnh để sinh câu trả lời tư vấn. Cách xử lý này giúp chatbot phản hồi dựa trên dữ liệu của đúng cửa hàng, thay vì trả lời chung chung hoặc sử dụng dữ liệu không thuộc phạm vi tư vấn.

Knowledge base được thiết kế như thành phần trung gian giữa dữ liệu cửa hàng và chatbot. Thay vì yêu cầu cửa hàng tự xây dựng hệ thống truy xuất phức tạp, hệ thống cung cấp chức năng quản lý nguồn dữ liệu và rebuild knowledge base. Khi cửa hàng cập nhật nguồn dữ liệu, hệ thống có thể xử lý lại dữ liệu để phục vụ truy xuất trong các phiên tư vấn sau. Nhờ đó, quy trình cập nhật tri thức cho chatbot được đưa về một nhóm thao tác quản trị rõ ràng, phù hợp hơn với người dùng phía cửa hàng.

Một điểm quan trọng trong giải pháp là liên kết giữa hội thoại và nghiệp vụ bán hàng. Trong quá trình tư vấn, nếu người dùng thể hiện ý định mua hàng hoặc cung cấp thông tin liên hệ, hệ thống có thể tạo lead hoặc yêu cầu mua hàng từ hội thoại. Các bản ghi này lưu lại thông tin người dùng, nhu cầu, sản phẩm quan tâm và nội dung hội thoại liên quan. Sau đó, nhân viên cửa hàng có thể truy cập danh sách lead hoặc yêu cầu mua hàng để claim, phản hồi, phân công hoặc cập nhật trạng thái xử lý. Nhờ vậy, chatbot không chỉ là công cụ trả lời tự động, mà còn trở thành điểm đầu vào cho quy trình bán hàng của cửa hàng.

Giải pháp cũng tách riêng backend nghiệp vụ và AI service. Backend chịu trách nhiệm xác thực, phân quyền, tenant isolation, quản lý dữ liệu nghiệp vụ và giao tiếp với frontend hoặc kênh ngoài. AI service tập trung vào xử lý hội thoại, truy xuất knowledge base và sinh câu trả lời. Cách tách này giúp hệ thống dễ bảo trì hơn, vì phần nghiệp vụ và phần AI có thể được phát triển, kiểm thử hoặc thay đổi tương đối độc lập. Khi cần thay đổi mô hình ngôn ngữ, chiến lược truy xuất hoặc cách xây dựng prompt, phần thay đổi chủ yếu nằm ở AI service mà không ảnh hưởng lớn đến các chức năng quản trị nghiệp vụ.

\subsection{Kết quả đạt được}
\label{subsection:5.1.3}

Kết quả đạt được của giải pháp là một thiết kế hệ thống có khả năng kết hợp ba thành phần chính: multi-tenant, RAG chatbot và quy trình xử lý lead/yêu cầu mua hàng. Về mặt tổ chức dữ liệu, hệ thống xác định tenant là đơn vị trung tâm để phân tách dữ liệu giữa các cửa hàng. Các chatbot, nguồn dữ liệu knowledge base, hội thoại, lead và yêu cầu mua hàng đều được gắn với tenant tương ứng. Cách tổ chức này giúp hệ thống phục vụ được nhiều cửa hàng mà không cần triển khai riêng một hệ thống cho từng cửa hàng.

Về mặt tư vấn, hệ thống sử dụng hướng tiếp cận RAG để chatbot có thể trả lời dựa trên dữ liệu riêng của cửa hàng. Khi người dùng hỏi về sản phẩm, hệ thống không chỉ gửi câu hỏi trực tiếp cho mô hình ngôn ngữ, mà còn truy xuất thông tin liên quan từ knowledge base để bổ sung ngữ cảnh trả lời. Điều này giúp chatbot có khả năng bám sát dữ liệu cửa hàng hơn so với một chatbot chỉ trả lời dựa trên kiến thức chung.

Về mặt nghiệp vụ, hệ thống không dừng lại ở hội thoại mà còn ghi nhận các thông tin có giá trị cho cửa hàng. Khi người dùng có ý định mua hàng hoặc để lại thông tin liên hệ, hệ thống có thể tạo lead hoặc yêu cầu mua hàng. Nhân viên cửa hàng có thể tiếp nhận các thông tin này, xem lại nội dung hội thoại và cập nhật trạng thái xử lý. Kết quả này giúp liên kết giữa quá trình tư vấn tự động và quy trình bán hàng thủ công của cửa hàng.

Về mặt kiến trúc, việc tách backend nghiệp vụ và AI service giúp hệ thống rõ ràng hơn về trách nhiệm của từng thành phần. Backend tập trung vào quản lý dữ liệu, phân quyền và nghiệp vụ multi-tenant; AI service tập trung vào xử lý hội thoại và truy xuất tri thức. Cách tổ chức này cũng tạo điều kiện để hệ thống có thể mở rộng trong tương lai, ví dụ bổ sung kênh chat ngoài, thay đổi mô hình ngôn ngữ, cải thiện thuật toán truy xuất hoặc mở rộng chức năng quản lý knowledge base.

\section{Giải pháp phân tách dữ liệu tenant và dữ liệu tham chiếu công khai}
\label{section:5.2}

\subsection{Bài toán đặt ra}
\label{subsection:5.2.1}

Trong quá trình phát triển hệ thống, một vấn đề quan trọng là làm thế nào để vừa đảm bảo nguyên tắc tách biệt dữ liệu tenant, vừa hỗ trợ các chức năng tư vấn tổng quát và tham chiếu giá. Đối với chatbot bán hàng theo cửa hàng, dữ liệu phải được xử lý nghiêm ngặt theo tenant. Chatbot của một cửa hàng không được lấy chính sách, sản phẩm hoặc hội thoại của cửa hàng khác để tư vấn cho khách hàng. Đây là yêu cầu cốt lõi của hệ thống multi-tenant.

Tuy nhiên, hệ thống cũng có các chức năng không đại diện cho một cửa hàng cụ thể, chẳng hạn tư vấn tổng quát và tham chiếu giá cơ bản. Các chức năng này cần dữ liệu rộng hơn phạm vi một cửa hàng để có thể đưa ra so sánh hoặc nhận xét giá. Nếu sử dụng trực tiếp dữ liệu riêng của nhiều tenant và gọi đó là dữ liệu tổng hợp, thiết kế sẽ dễ gây hiểu nhầm rằng dữ liệu riêng của các cửa hàng bị trộn với nhau. Điều này không phù hợp với nguyên tắc multi-tenant và cũng gây khó khăn khi giải thích thiết kế hệ thống.

Bài toán đặt ra là cần thiết kế một cơ chế phân loại dữ liệu rõ ràng. Dữ liệu riêng của cửa hàng phải tiếp tục được cô lập theo tenant, trong khi các chức năng tư vấn chung chỉ được sử dụng nguồn dữ liệu có phạm vi công khai hoặc tham chiếu. Cách tổ chức dữ liệu cần đủ rõ để tránh việc chat theo cửa hàng sử dụng nhầm dữ liệu tham chiếu, hoặc ngược lại, chat tham chiếu sử dụng nhầm dữ liệu riêng của một tenant.

\subsection{Giải pháp đề xuất}
\label{subsection:5.2.2}

Giải pháp được đề xuất là phân tách dữ liệu thành hai phạm vi chính: dữ liệu theo tenant và dữ liệu tham chiếu công khai. Dữ liệu theo tenant là dữ liệu phục vụ trực tiếp cho một cửa hàng, bao gồm sản phẩm, chính sách, bài viết hoặc nguồn tri thức được cửa hàng cấu hình. Dữ liệu này phải gắn với \texttt{tenant\_id} và chỉ được sử dụng trong các luồng thuộc tenant tương ứng.

Dữ liệu tham chiếu công khai là dữ liệu được thu thập từ các nguồn công khai, phục vụ cho các chức năng không đại diện cho một cửa hàng cụ thể như \texttt{general\_compare} và \texttt{market\_price}. Dữ liệu này không được hiểu là dữ liệu riêng của tenant bị gom lại, mà là một kho tham chiếu độc lập. Khi sử dụng dữ liệu tham chiếu công khai, hệ thống cần giữ lại nguồn gốc dữ liệu như tên cửa hàng, đường dẫn nguồn và thời điểm thu thập để người dùng có thể hiểu phạm vi tham khảo của kết quả.

Trong pipeline dữ liệu, mỗi nguồn dữ liệu được gắn với trường \texttt{visibility}. Trường này có thể nhận các giá trị như \texttt{tenant\_private}, \texttt{tenant\_public} hoặc \texttt{public\_reference}. Dữ liệu \texttt{tenant\_private} và \texttt{tenant\_public} phải có \texttt{tenant\_id}; dữ liệu \texttt{public\_reference} không mang \texttt{tenant\_id} như dữ liệu nghiệp vụ riêng của một cửa hàng. Cách này giúp hệ thống tránh nhầm lẫn giữa dữ liệu riêng và dữ liệu tham chiếu.

Quy tắc sử dụng dữ liệu theo mode được xác định như sau. Chế độ \texttt{tenant\_sales} chỉ sử dụng dữ liệu có \texttt{tenant\_id} khớp với tenant đang được phục vụ. Chế độ \texttt{general\_compare} chỉ sử dụng dữ liệu \texttt{public\_reference}. Chế độ \texttt{market\_price} cũng chỉ sử dụng dữ liệu \texttt{public\_reference} để đưa ra khoảng giá tham chiếu trong phạm vi dữ liệu hệ thống. Các dữ liệu chính sách hoặc giới thiệu cửa hàng chỉ được sử dụng trong phạm vi tenant tương ứng.

Ở mức tổ chức file, dữ liệu sau xử lý được chia thành các thư mục theo phạm vi sử dụng. Dữ liệu riêng của tenant được đưa vào thư mục \texttt{output/tenants/\{tenant\_id\}}, còn dữ liệu tham chiếu được đưa vào \texttt{output/reference}. Việc tách thư mục giúp dễ kiểm tra dữ liệu đầu ra và hỗ trợ validate trước khi đưa vào runtime chatbot.

\subsection{Kết quả đạt được}
\label{subsection:5.2.3}

Kết quả của giải pháp là hệ thống có được một cách tổ chức dữ liệu rõ ràng hơn cho môi trường multi-tenant. Dữ liệu bán hàng theo cửa hàng vẫn được cô lập theo tenant, trong khi các chức năng tư vấn tổng quát và tham chiếu giá sử dụng một kho tham chiếu công khai riêng. Cách thiết kế này giúp tránh hiểu nhầm rằng hệ thống trộn dữ liệu riêng của nhiều cửa hàng.

Về mặt triển khai, pipeline dữ liệu có thể kiểm tra được các lỗi phổ biến như dữ liệu tenant thiếu \texttt{tenant\_id}, dữ liệu tham chiếu công khai lại mang \texttt{tenant\_id}, hoặc record được ghi sai thư mục đầu ra. Những kiểm tra này giúp phát hiện sớm lỗi dữ liệu trước khi tích hợp vào chatbot runtime.

Về mặt trình bày sản phẩm, giải pháp cũng giúp hệ thống dễ giải thích hơn. Khi demo hoặc báo cáo, có thể nêu rõ rằng chat bán hàng theo cửa hàng chỉ dùng dữ liệu tenant tương ứng, còn chat chung và khảo giá dùng dữ liệu tham chiếu công khai. Đây là điểm quan trọng để bảo đảm hệ thống vừa đáp ứng yêu cầu multi-tenant, vừa hỗ trợ được các chức năng tham khảo cho người dùng cuối.

% TODO_CH5_DATA_PIPELINE_RESULT: Nếu đã crawl dữ liệu thật, bổ sung số lượng nguồn, số sản phẩm và số knowledge document đã thu thập.
% Không tự điền nếu chưa chạy validate_catalog.py.

\section{Giải pháp kiểm soát ba chế độ hội thoại}
\label{section:5.3}

\subsection{Bài toán đặt ra}
\label{subsection:5.3.1}

Trong hệ thống chatbot nội thất, không phải mọi hội thoại đều có cùng mục tiêu. Một người dùng có thể chat với chatbot của một cửa hàng cụ thể để hỏi mua sản phẩm, nhưng cũng có thể chỉ muốn hỏi tư vấn chung hoặc tham khảo mức giá. Nếu tất cả các luồng chat được xử lý giống nhau, hệ thống có thể tạo ra các hành vi không phù hợp. Ví dụ, một cuộc hội thoại tham chiếu giá có thể vô tình tạo yêu cầu mua hàng, hoặc một câu hỏi so sánh chung có thể bị xử lý như luồng bán hàng của một tenant.

Vấn đề này trở nên quan trọng hơn vì hệ thống có liên kết giữa hội thoại và nghiệp vụ bán hàng. Việc tạo lead hoặc purchase request là thao tác ghi dữ liệu nghiệp vụ, ảnh hưởng đến quy trình xử lý của cửa hàng. Vì vậy, hệ thống cần kiểm soát chặt chẽ khi nào được phép tạo yêu cầu mua hàng, khi nào chỉ được trả lời tham khảo.

Bài toán đặt ra là cần xác định rõ các chế độ hội thoại và contract giữa backend với AI service. Mỗi chế độ cần có phạm vi dữ liệu, hành vi trả lời và tác động nghiệp vụ khác nhau. Hệ thống cũng cần có cơ chế bảo vệ ở backend, không phụ thuộc hoàn toàn vào phản hồi của mô hình ngôn ngữ.

\subsection{Giải pháp đề xuất}
\label{subsection:5.3.2}

Giải pháp được đề xuất là chuẩn hóa ba chế độ hội thoại chính: \texttt{tenant\_sales}, \texttt{general\_compare} và \texttt{market\_price}. Mỗi chế độ có mục tiêu và phạm vi xử lý riêng.

Chế độ \texttt{tenant\_sales} phục vụ tư vấn bán hàng theo một cửa hàng cụ thể. Chế độ này sử dụng dữ liệu của tenant tương ứng, có thể hỏi làm rõ nhu cầu, tư vấn sản phẩm và tạo yêu cầu mua hàng khi người dùng cung cấp đủ thông tin cần thiết. Đây là chế độ duy nhất được phép phát sinh lead hoặc purchase request.

Chế độ \texttt{general\_compare} phục vụ tư vấn và so sánh chung. Chế độ này không đại diện cho một cửa hàng cụ thể và sử dụng dữ liệu tham chiếu công khai nếu có. Mục tiêu của chế độ này là giúp người dùng hiểu các lựa chọn sản phẩm, tiêu chí so sánh và ưu nhược điểm ở mức tham khảo. Chế độ này không tạo lead hoặc purchase request.

Chế độ \texttt{market\_price} phục vụ tham chiếu giá cơ bản. Chế độ này dùng dữ liệu tham chiếu công khai để đưa ra khoảng giá hoặc nhận xét một mức giá cụ thể trong phạm vi dữ liệu hệ thống. Chế độ này không thay thế cho việc định giá thị trường đầy đủ và không đề xuất mua hàng cụ thể. Nếu dữ liệu không đủ, hệ thống cần trả lời rõ rằng chưa đủ nguồn tham chiếu, thay vì để mô hình tự ước lượng giá.

Trong contract giữa backend và AI service, mode được gửi kèm trong request. AI service dùng mode để chọn prompt, chiến lược truy xuất dữ liệu và logic phản hồi phù hợp. Kết quả trả về có thể bao gồm mode cuối cùng, debug trace và cờ cho biết có đủ điều kiện tạo purchase request hay không. Tuy nhiên, backend vẫn là lớp kiểm soát cuối cùng. Backend chỉ tạo purchase request khi mode yêu cầu là \texttt{tenant\_sales}, mode cuối cùng do AI service trả về cũng là \texttt{tenant\_sales}, conversation thuộc tenant hợp lệ và các điều kiện nghiệp vụ khác được thỏa mãn.

Cách thiết kế này tạo ra hai lớp bảo vệ. Lớp thứ nhất nằm ở AI service, nơi prompt và logic hội thoại được điều chỉnh theo mode. Lớp thứ hai nằm ở backend, nơi các tác động nghiệp vụ như tạo purchase request được kiểm tra lại. Nhờ đó, ngay cả khi AI service trả về tín hiệu không phù hợp trong một luồng không bán hàng, backend vẫn có thể chặn tác động nghiệp vụ sai.

\subsection{Kết quả đạt được}
\label{subsection:5.3.3}

Kết quả của giải pháp là hệ thống phân biệt rõ ba nhóm nhu cầu hội thoại: tư vấn bán hàng theo cửa hàng, tư vấn/so sánh chung và tham chiếu giá. Việc chuẩn hóa mode giúp giảm tình trạng các luồng chat bị lẫn hành vi với nhau. Đặc biệt, hệ thống có thể bảo đảm rằng \texttt{general\_compare} và \texttt{market\_price} không tạo purchase request, trong khi \texttt{tenant\_sales} vẫn giữ khả năng chuyển hội thoại thành yêu cầu mua hàng khi phù hợp.

Giải pháp này cũng giúp quá trình kiểm thử trở nên rõ ràng hơn. Các test case có thể kiểm tra riêng từng mode, bao gồm dữ liệu được sử dụng, phản hồi của chatbot và tác động nghiệp vụ sau hội thoại. Ví dụ, một test case cho \texttt{market\_price} cần kiểm tra rằng hệ thống trả lời trong phạm vi dữ liệu tham chiếu và không tạo purchase request. Một test case cho \texttt{tenant\_sales} cần kiểm tra rằng purchase request chỉ được tạo khi hội thoại đủ điều kiện.

Về mặt vận hành, contract mode giúp backend và AI service phối hợp rõ ràng hơn. Khi cần bổ sung mode mới trong tương lai, hệ thống có thể mở rộng theo cùng nguyên tắc: xác định phạm vi dữ liệu, hành vi trả lời, tác động nghiệp vụ được phép và các điều kiện kiểm soát ở backend.

% TODO_CH5_MODE_TEST_RESULT: Nếu có log/test đã verify 3 mode, chèn tóm tắt kết quả vào đây.
% Ví dụ: số test pass, log chặn purchase request ở general_compare/market_price, hoặc screenshot debug mode.

\section{Giải pháp pipeline dữ liệu cho catalog sản phẩm và knowledge base}
\label{section:5.4}

\subsection{Bài toán đặt ra}
\label{subsection:5.4.1}

Chất lượng câu trả lời của chatbot phụ thuộc lớn vào chất lượng dữ liệu đầu vào. Với bài toán nội thất, dữ liệu cần không chỉ có tên sản phẩm mà còn có giá, chất liệu, kích thước, loại sản phẩm, phong cách, mô tả, cửa hàng và đường dẫn nguồn. Ngoài dữ liệu sản phẩm, chatbot theo cửa hàng còn cần các thông tin như chính sách bảo hành, vận chuyển, giới thiệu cửa hàng và bài viết hướng dẫn chọn sản phẩm.

Nếu nhập dữ liệu thủ công hoàn toàn, việc chuẩn bị dữ liệu sẽ mất nhiều thời gian và khó mở rộng. Nếu crawler tự động quét toàn bộ website, dữ liệu dễ bị nhiễu, lấy nhầm link hoặc thiếu kiểm soát. Ngoài ra, nếu dữ liệu crawl được đưa thẳng vào chatbot mà không có bước validate, chatbot có thể trả lời dựa trên dữ liệu thiếu, trùng lặp hoặc sai phạm vi tenant.

Bài toán đặt ra là cần một pipeline dữ liệu có kiểm soát. Pipeline này phải cho phép người phát triển chọn nguồn dữ liệu, crawl dữ liệu, chuẩn hóa thành catalog có cấu trúc, validate chất lượng và tách dữ liệu theo phạm vi sử dụng trước khi đưa vào chatbot.

\subsection{Giải pháp đề xuất}
\label{subsection:5.4.2}

Giải pháp được đề xuất là xây dựng một thư mục \texttt{data\_pipeline} độc lập với runtime chatbot. Thư mục này chứa các file CSV đầu vào, script crawl, parser, script normalize và script validate. Cách tách riêng này giúp quá trình chuẩn bị dữ liệu không ảnh hưởng trực tiếp đến service đang chạy, đồng thời giúp dễ kiểm tra và tái chạy pipeline khi cần cập nhật dữ liệu.

Nguồn dữ liệu đầu vào được quản lý bằng file CSV. Mỗi dòng trong CSV mô tả một nguồn cần xử lý, gồm \texttt{tenant\_id}, tên cửa hàng, loại nguồn, danh mục, mức visibility, URL và ghi chú. Loại nguồn có thể là sản phẩm cụ thể, danh mục sản phẩm, chính sách hoặc bài viết nội dung. Cách quản lý bằng CSV giúp người phát triển kiểm soát rõ crawler sẽ truy cập những URL nào, tránh việc crawler tự đi quét toàn bộ website một cách khó kiểm soát.

Dữ liệu sau crawl được chia thành hai nhóm đầu ra. Nhóm thứ nhất là \texttt{products.clean.jsonl}, chứa các bản ghi sản phẩm có cấu trúc như tên, loại, giá, đơn vị tiền tệ, chất liệu, kích thước, phong cách, mô tả, ảnh, URL nguồn và thời điểm crawl. Nhóm thứ hai là \texttt{knowledge\_docs.clean.jsonl}, chứa các tài liệu dạng chính sách, giới thiệu hoặc bài viết, gồm tiêu đề, nội dung, loại tài liệu, URL nguồn và thời điểm crawl.

Pipeline cũng có bước validate để kiểm tra chất lượng dữ liệu. Với sản phẩm, validate cần thống kê số lượng sản phẩm, số lượng theo nguồn, tỷ lệ có giá, tỷ lệ có chất liệu/kích thước, số URL trùng và các trường thiếu. Với tài liệu knowledge, validate cần kiểm tra tiêu đề, nội dung, nguồn và tenant nếu có. Các lỗi về visibility và tenant cũng cần được phát hiện, ví dụ dữ liệu tenant thiếu \texttt{tenant\_id} hoặc dữ liệu public reference lại bị ghi vào thư mục tenant.

Nguyên tắc quan trọng của pipeline là không bịa dữ liệu. Nếu website không cung cấp giá, chất liệu hoặc kích thước, field tương ứng được để \texttt{null}. Nếu cần dùng mô hình ngôn ngữ để hỗ trợ chuẩn hóa mô tả trong tương lai, pipeline vẫn phải giữ lại raw text và source URL để truy vết, đồng thời không được tự tạo ra trường dữ liệu không có căn cứ.

\subsection{Kết quả đạt được}
\label{subsection:5.4.3}

Kết quả của giải pháp là hệ thống có một hướng chuẩn bị dữ liệu rõ ràng hơn so với việc nhập tay rời rạc hoặc dùng dữ liệu mock. Dữ liệu sản phẩm và tài liệu cửa hàng được đưa vào một quy trình có các bước rõ ràng: khai báo nguồn, crawl, parse, normalize, validate và xuất dữ liệu sạch.

Về mặt kiểm soát chất lượng, pipeline giúp phát hiện dữ liệu thiếu, dữ liệu trùng hoặc dữ liệu sai phạm vi trước khi tích hợp vào chatbot. Điều này đặc biệt quan trọng đối với hai chức năng \texttt{general\_compare} và \texttt{market\_price}, vì hai chức năng này cần dữ liệu có cấu trúc để so sánh sản phẩm hoặc tính khoảng giá tham chiếu. Nếu dữ liệu không đủ, hệ thống có thể trả lời rằng chưa đủ nguồn tham chiếu thay vì đưa ra kết luận không có căn cứ.

Về mặt phát triển, việc tách \texttt{data\_pipeline} khỏi runtime chatbot giúp quá trình thử nghiệm an toàn hơn. Người phát triển có thể crawl thử với số lượng URL nhỏ, kiểm tra report, chỉnh parser, sau đó mới mở rộng dữ liệu. Khi dữ liệu đạt yêu cầu, file output sạch mới được copy hoặc cấu hình cho chatbot sử dụng.

% TODO_CH5_CRAWL_RESULT: Chèn kết quả validate_catalog.py khi đã crawl thật.
% Cần có: tổng số sản phẩm, số store, tỷ lệ có price/material/size, số knowledge docs, số lỗi duplicate/missing.

\section{Giải pháp triển khai nhẹ với Claude API và local model fallback}
\label{section:5.5}

\subsection{Bài toán đặt ra}
\label{subsection:5.5.1}

Một khó khăn thực tế trong quá trình triển khai chatbot dùng mô hình ngôn ngữ là yêu cầu tài nguyên. Nếu chạy local model trên CPU, thời gian phản hồi có thể dài và bộ nhớ sử dụng có thể tăng cao. Nếu sử dụng GPU, chi phí thuê máy chủ sẽ lớn hơn đáng kể. Trong khi đó, mục tiêu của đồ án là triển khai được một hệ thống demo ổn định trên VPS phổ thông, có khả năng phục vụ các luồng chính như chat theo cửa hàng, chat tổng quát và khảo giá.

Phiếu giao nhiệm vụ có đề cập đến việc sử dụng mô hình local và API bên ngoài để so sánh hoặc thay thế trong các trường hợp khác nhau. Tuy nhiên, điều này không đồng nghĩa với việc hệ thống triển khai cuối cùng bắt buộc phải chạy local model liên tục trên VPS. Nếu cố chạy local model trong môi trường CPU-only cấu hình thấp, hệ thống có thể trở nên chậm, khó kiểm soát và ảnh hưởng đến trải nghiệm demo.

Bài toán đặt ra là cần thiết kế cách triển khai cân bằng giữa chất lượng phản hồi, chi phí vận hành và khả năng chạy ổn định trên hạ tầng hạn chế. Hệ thống cần dùng provider chất lượng cao cho luồng chính, đồng thời vẫn giữ khả năng thử nghiệm hoặc fallback local model khi được bật rõ ràng.

\subsection{Giải pháp đề xuất}
\label{subsection:5.5.2}

Giải pháp được lựa chọn là sử dụng Claude API làm provider chính trong môi trường triển khai, còn Qwen local được giữ như một phương án thử nghiệm hoặc fallback kỹ thuật. API key của Claude được cấu hình bằng biến môi trường như \texttt{ANTHROPIC\_API\_KEY} hoặc \texttt{CLAUDE\_API\_KEY}, không lưu trong mã nguồn, tài liệu hoặc cơ sở dữ liệu. Cách cấu hình này phù hợp với yêu cầu bảo vệ khóa tích hợp và giúp dễ thay đổi key khi cần.

Trong Docker Compose, hệ thống gồm ba service chính: \texttt{app}, \texttt{chatbot-api} và \texttt{postgres}. Backend Spring Boot chạy trong service \texttt{app}; Python FastAPI chạy trong service \texttt{chatbot-api}; PostgreSQL chạy trong service \texttt{postgres}. Backend gọi AI service thông qua \texttt{PYTHON\_LLM\_BASE\_URL=http://chatbot-api:8000}. Cách tổ chức này giúp Spring Boot không cần tự spawn Python process trong môi trường triển khai, giảm phụ thuộc vào đường dẫn Python cục bộ và làm cho hệ thống dễ chạy lại trên VPS hơn.

Để tránh local model tiêu tốn tài nguyên ngoài ý muốn, cấu hình triển khai mặc định đặt \texttt{FALLBACK\_TO\_LOCAL\_ENABLED=false}. Khi biến này tắt và Claude API được cấu hình, Python service không warmup hoặc load Qwen khi khởi động. Local pipeline chỉ được load khi provider local được bật rõ ràng. Ngoài ra, hệ thống bổ sung cấu hình giới hạn cache local pipeline và cơ chế unload sau một khoảng thời gian không sử dụng. Cơ chế này giúp giảm rủi ro bộ nhớ tăng cao nếu local model được bật trong quá trình thử nghiệm.

Với VPS, hệ thống sử dụng file \texttt{.env} để lưu cấu hình triển khai. File này không được commit vào Git và được giới hạn quyền đọc trên server. Các port nội bộ như PostgreSQL và chatbot-api không cần mở public; người dùng truy cập chủ yếu thông qua service backend hoặc reverse proxy. Khi cần webhook Messenger/Telegram ổn định, hệ thống có thể dùng domain và HTTPS thông qua Caddy hoặc Nginx.

\subsection{Kết quả đạt được}
\label{subsection:5.5.3}

Kết quả của giải pháp là hệ thống có thể chạy theo hướng nhẹ hơn trên VPS CPU-only. Các service chính được đóng gói bằng Docker Compose, backend gọi Python service qua HTTP nội bộ và provider chính là Claude API. Do local model không được load mặc định, yêu cầu RAM và CPU khi khởi động giảm đáng kể so với hướng chạy local model liên tục.

Cách triển khai này cũng giúp giảm rủi ro phụ thuộc vào môi trường local của máy phát triển. Trước đó, khi chạy bằng môi trường cục bộ, hệ thống có thể phụ thuộc vào đường dẫn Python, thư mục model hoặc biến môi trường riêng của máy cá nhân. Khi chuyển sang Docker Compose, các service được cấu hình rõ ràng hơn, giúp việc chạy lại trên VPS hoặc máy khác dễ kiểm soát hơn.

Về mặt bảo mật, việc chuyển API key sang biến môi trường và loại bỏ cấu hình key theo từng chatbot giúp giảm nguy cơ lộ key qua database hoặc giao diện quản trị. Nếu cần thay đổi hoặc vô hiệu hóa key, người vận hành chỉ cần chỉnh biến môi trường hoặc thao tác trên console của provider.

% TODO_CH5_DEPLOY_RESULT: Chèn kết quả deploy thật nếu muốn.
% Cần có: docker compose ps, free -h, healthz, log warmup skip local model, URL truy cập hệ thống.

\section{Tổng kết đóng góp}
\label{section:5.6}

Các giải pháp trình bày trong chương này tập trung vào những vấn đề cốt lõi của đồ án. Thứ nhất, hệ thống kết hợp multi-tenant với RAG chatbot để phục vụ tư vấn bán hàng nội thất cho nhiều cửa hàng trên cùng một nền tảng. Thứ hai, hệ thống phân biệt dữ liệu tenant và dữ liệu tham chiếu công khai để vừa đảm bảo tenant isolation, vừa hỗ trợ được chat tổng quát và tham chiếu giá. Thứ ba, hệ thống chuẩn hóa ba chế độ hội thoại nhằm kiểm soát rõ phạm vi dữ liệu và tác động nghiệp vụ của từng luồng. Thứ tư, hệ thống xây dựng pipeline dữ liệu riêng để chuẩn bị catalog sản phẩm và knowledge base một cách có kiểm soát. Cuối cùng, hệ thống được triển khai theo hướng nhẹ bằng Docker Compose và Claude API, giữ local model như phương án thử nghiệm hoặc fallback thay vì phụ thuộc vào GPU.

Những đóng góp này không nhằm xây dựng một nền tảng AI agent tổng quát, mà tập trung giải quyết một bài toán cụ thể: hỗ trợ tư vấn bán hàng nội thất trong môi trường nhiều cửa hàng, có dữ liệu riêng, có hội thoại nhiều lượt và có khả năng chuyển kết quả tư vấn thành thông tin nghiệp vụ để cửa hàng tiếp tục xử lý. Đây là điểm khác biệt chính của đồ án so với một chatbot hỏi đáp thông thường.

\end{document}